\documentclass[10pt,a4paper]{amsart}
\usepackage[margin=19mm]{geometry}
\usepackage{amsmath,amssymb,mathtools}
\usepackage[scr=rsfs,cal=euler]{mathalfa}
\usepackage{xfrac}
\usepackage[unicode,hidelinks,bookmarks=false]{hyperref}
\newcommand{\ie}{i.e.\ }
\newcommand{\cf}{cf.\ }
\newcommand{\resp}{respectively\ }
\newcommand{\Subsection}{\subsection}
\providecommand{\nobreakdash}{\nobreak\hbox{-}}
\setlength{\emergencystretch}{2em}
\multlinegap0pt
\usepackage{amssymb}
\usepackage{mathtools}
\usepackage{bm}
\usepackage{xcolor}
\usepackage{float}
\usepackage[shortlabels]{enumitem}
\newtheorem{assumption}{Assumption}
\newtheorem{lemma}{Lemma}
\title{Quenched Amplification and Tail Shaping in Networked Systems with Memory and Regime Switching}
\author{Mauricio Herrera-Mar\'in}
\date{Source: arXiv:2605.00750v1 (CC BY 4.0)}
\begin{document}
\maketitle

\noindent\emph{Source and licence.} Quenched Amplification and Tail Shaping in Networked Systems with Memory and Regime Switching. Version arXiv:2605.00750v1 (CC BY 4.0), distributed by arXiv under the Creative Commons Attribution 4.0 International licence, http://creativecommons.org/licenses/by/4.0/, as recorded in the arXiv licence metadata for this exact version and shown on the abstract page https://arxiv.org/abs/2605.00750v1. Full licence notice: \texttt{licenses/arxiv-2605.00750-CC-BY-4.0.txt}.\par\smallskip

\noindent\textit{Editorial scope.} Counted unit: the article's lemma lem:cone\_projected\_growth (source0.tex lines 425--436). The article's complete original proof of it is delivered with it (source0.tex lines 438--480, one \texttt{proof} environment of 43 lines). No mathematics is written, repaired or re-derived: the statement and the proof are the article's own lines, in the article's own notation, and no retained line is substituted or re-flowed.  Retained uncounted as the source-local context the proof needs: lines 379--383; lines 387--391; lines 398--405; lines 407--413.  Nothing was omitted from any retained span, so the package carries no omission record.  No retained line contains a citation, so the wrapper carries no bibliography and no bibliography entry.  No retained line contains a figure environment, an image-inclusion command or drawing code, so no image is delivered and the package declares no asset dependency.  Editorial formatting only, in addition: an AMS \texttt{amsart} preamble that restates the article's own declarations and macros used by the retained lines, namely the environments \texttt{assumption}, \texttt{lemma} on separate counters, together with the standard packages those lines need.  The retained environments print in this document as Assumption 1, Lemma 1 in order of appearance, whereas the article prints the same items with the numbering of its own full document.  The complete native source of the article as distributed in the arXiv e-print for this exact version is bundled unchanged as \texttt{sources/199589/source0.tex}.
\begin{equation}\label{eq:vU_top_def}
	\|v_U\|_2=1,
	\qquad
	v_U^\top A_U^{(0)} v_U=\gamma_U.
\end{equation}

\begin{assumption}[Unfavorable instantaneous susceptibility]\label{ass:gammaU_positive_cone}
	\begin{equation}\label{eq:gammaU_pos_cone}
		\gamma_U>0.
	\end{equation}
\end{assumption}

\begin{assumption}[Cone condition on amplifying $U$-dwells]\label{ass:cone_condition}
	There exist $\alpha\in(0,1]$ and $T_0\in(0,T]$ such that, on every sample path segment where an extreme burst is produced,
	there exists a sub-interval $[t_0,t_0+\tau]\subset[0,T_0]$ with $z(t)\equiv U$ and the associated solution satisfies
	\begin{equation}\label{eq:cone}
		v_U^\top X(t)\ \ge\ \alpha\,\|X(t)\|_2
		\qquad \text{for all } t\in[t_0,t_0+\tau].
	\end{equation}
\end{assumption}

\begin{assumption}[Persistent forcing injection along the cone]\label{ass:forcing_in_cone}
	There exists $c_f>0$ such that
	\begin{equation}\label{eq:forcing_cone}
		v_U^\top \widetilde f(t)\ \ge\ c_f
		\qquad \text{for a.e.\ } t\in[0,T_0].
	\end{equation}
\end{assumption}

\begin{lemma}[Cone-projected growth inequality]\label{lem:cone_projected_growth}
	Assume \ref{ass:gammaU_positive_cone}, \ref{ass:cone_condition}, and \ref{ass:forcing_in_cone}.
	Suppose that for some $t_0\ge 0$ and $\tau>0$ with $t_0+\tau\le T_0$ we have $z(t)\equiv U$ and $m(t)\equiv 0$ on $[t_0,t_0+\tau]$.
	Define $\zeta(t):=v_U^\top X(t)$. Then for almost every $t\in[t_0,t_0+\tau]$,
	\begin{equation}\label{eq:zeta_cone_ineq}
		\dot\zeta(t)\ \ge\ \gamma_U\,\zeta(t)+c_f,
	\end{equation}
	and hence
	\begin{equation}\label{eq:zeta_cone_solution}
		\zeta(t_0+\tau)\ \ge\ e^{\gamma_U\tau}\,\zeta(t_0)+\frac{c_f}{\gamma_U}\Big(e^{\gamma_U\tau}-1\Big).
	\end{equation}
\end{lemma}

\begin{proof}
	On $[t_0,t_0+\tau]$, $\dot X=A_U^{(0)}X+\widetilde f$. Thus
	\[
	\dot\zeta(t)=v_U^\top A_U^{(0)}X(t)+v_U^\top\widetilde f(t).
	\]
	Write $X(t)=\zeta(t)v_U+w(t)$ with $w(t)\perp v_U$. Then
	\[
	v_U^\top A_U^{(0)}X(t)
	=\zeta(t)\,v_U^\top A_U^{(0)}v_U + v_U^\top A_U^{(0)}w(t)
	=\gamma_U\,\zeta(t) + v_U^\top A_U^{(0)}w(t),
	\]
	using \eqref{eq:vU_top_def}. The cone condition \eqref{eq:cone} implies
	$\zeta(t)=v_U^\top X(t)\ge \alpha\|X(t)\|_2$, hence $\|w(t)\|_2^2=\|X(t)\|_2^2-\zeta(t)^2\le (1-\alpha^2)\|X(t)\|_2^2$ and therefore
	\begin{equation}\label{eq:w_bound}
		\|w(t)\|_2 \le \sqrt{1-\alpha^2}\,\|X(t)\|_2 \le \frac{\sqrt{1-\alpha^2}}{\alpha}\,\zeta(t)
		\qquad \text{for all } t\in[t_0,t_0+\tau].
	\end{equation}
	Consequently,
	\[
	v_U^\top A_U^{(0)}w(t)\ \ge\ -\|v_U^\top A_U^{(0)}\|_2\,\|w(t)\|_2
	\ge
	-\|A_U^{(0)}\|_2\,\frac{\sqrt{1-\alpha^2}}{\alpha}\,\zeta(t).
	\]
	Define the effective cone growth rate
	\begin{equation}\label{eq:gammaU_eff}
		\gamma_U^{\mathrm{cone}}
		:=
		\gamma_U - \|A_U^{(0)}\|_2\,\frac{\sqrt{1-\alpha^2}}{\alpha}.
	\end{equation}
	Then
	\[
	v_U^\top A_U^{(0)}X(t)
	\ge
	\gamma_U^{\mathrm{cone}}\,\zeta(t).
	\]
	Combining with \eqref{eq:forcing_cone} gives for a.e.\ $t$:
	\[
	\dot\zeta(t)\ge \gamma_U^{\mathrm{cone}}\zeta(t)+c_f.
	\]
	If $\gamma_U^{\mathrm{cone}}>0$, set $\gamma_U:=\gamma_U^{\mathrm{cone}}$ (redefining the effective growth rate on the cone),
	and the differential inequality becomes \eqref{eq:zeta_cone_ineq}. Solving it with the integrating factor
	$e^{-\gamma_U t}$ yields \eqref{eq:zeta_cone_solution}.
\end{proof}

\end{document}
